\section{Related Work}
\label{sec:related}

Existing simulators address important parts of this problem: multi-modality rendering, fluorescence-cell simulation, bacterial-cell synthetic training data, point-emitter benchmarks, and high-fidelity interference-microscopy simulation. DeepTrack and DeepTrack~2 \cite{deeptrack} provide a compositional framework for bright-field, dark-field, interferometric-scattering, holography, and fluorescence simulations, with Mie scatterers and flexible sample/optics pipelines. CytoPacq \cite{cytopacq} targets fluorescence-cell imaging with acquisition-artifact modeling, and SyMBac \cite{symbac} uses synthetic micrographs to train bacterial-cell segmentation models. Single-molecule localization microscopy benchmark generators provide established point-emitter simulations \cite{smlm-simulator}. Hitzelhammer et al. \cite{hitzelhammer-interference-platform} present a unified interference-microscopy simulation platform for interferometric scattering microscopy, coherent bright-field microscopy, and dark-field microscopy, with boundary-element scattered fields and Richards--Wolf vectorial imaging transformations. That work is closer to Syniscopy's label-free optical paths than general synthetic-data simulators and uses explicit scalar/vectorial coherent-backend choices aligned with Syniscopy's configurable optical path set. Syniscopy is organized around a different object: a shared single-particle scene whose modality-specific renders, detector-domain derivatives, noise maps, Fisher matrices, Cram\'{e}r--Rao summaries, fusion bounds, serial-allocation diagnostics, and supervision masks are defined on one physical parameter frame.

The mathematical components are established. Lateral localization bounds trace to Thompson, Larson and Webb \cite{thompson-larson-webb} and the Fisher information treatment of Ober, Ram and Ward \cite{ober-ram-ward}. Dong, Maestre, Conrad-Billroth and Juffmann \cite{juffmann-iscat-cobri-crlb} analyzed fundamental bounds for interferometric scattering microscopy, coherent bright-field imaging, and dark-field microscopy in a shared theoretical framework; Syniscopy includes those coherent limits while its bright-field and dark-field modality-comparison rows use experimentally oriented primary paths. Syniscopy also includes phase-contrast, fluorescence, transmission-electron-microscopy, and scanning-electron-microscopy response models. Fisher information additivity is standard in sensor fusion \cite{fisher-fusion}, and A-/D-/E-optimal allocation is standard in optimal design and convex optimization \cite{pukelsheim-optimal-design,boyd-convex}. Syniscopy integrates these established components into one simulator with a shared parameter frame; the new formal step is the contrast-functional symmetry theorem and fusion corollary on that shared frame, not a new Cram\'{e}r--Rao or optimization formalism.

Rigid-body localization and orientation precision have also been treated directly. McGray, Copeland, Stavis and Geist \cite{orientation-precision} extended localization precision from optical point sources to microscopic rigid bodies and derived centroid and orientation precision for planar localization microscopy. Syniscopy's $\mathrm{SE}(3)$ result is not a first rigid-body precision formula; it links modality-specific contrast-functional stabilizers to singular axes in rendered Fisher matrices, then reports the predicted rank deficits and fusion-recovered axes as simulator diagnostics.

Foundation segmentation models such as Segment Anything Model, Segment Anything Model 2 (SAM2) \cite{sam,sam2}, MedSAM \cite{medsam}, and micro-sam \cite{microsam} have made promptable segmentation and fine-tuning practical, but annotated microscopy video remains a bottleneck. Synthetic generation provides latent object geometry, yet particle videos also need a rule for pixels where the object is known but the rendered measurement does not support a reliable positive label. Loss weighting and boundary- or imbalance-aware segmentation losses are established training devices \cite{dice-loss,boundary-loss}; Syniscopy computes measurement-support scores from signal, Cram\'{e}r--Rao, temporal, and assignment-ambiguity factors, then provides geometry/support, supported, ignore, and loss-weight tensors that Segment Anything Model 2 or other segmentation models can use during training.

The label-free coherent-optical modes in Syniscopy include both scalar and vectorial paths. Fluorescence, transmission electron microscopy, and scanning electron microscopy route through the same simulator interface with explicit backend selectors and fidelity metadata, supporting both proxy diagnostics and higher-fidelity backends where configured. Learned generators such as generative adversarial networks for microscopy and diffusion models for biomedical images \cite{gan-microscopy,diffusion-microscopy} can produce visually realistic microscopy images, but they do not by themselves emit the latent particle state. They also do not provide the geometric masks, Fisher information, or measurement-supported ignore regions needed for this paper's cross-modality bound comparison. A Fisher bound computed on a generated image is a bound for that image; without a latent physical particle, calibrated detector model, and differentiable parameter frame, it is not a cross-modality instrument bound.
